\documentclass[12pt,a4paper,UTF8]{ctexart}
\usepackage[a4paper,top=22mm,bottom=22mm,left=22mm,right=22mm,headheight=15pt]{geometry}
\usepackage{amsmath,amssymb,mathtools,bm}
\usepackage{graphicx,xcolor,booktabs,longtable,array,enumitem,fancyhdr,caption,float,microtype,setspace}
\newcommand{\Needspace}[1]{\par\begingroup\dimen0=#1\relax\dimen2=\pagegoal\advance\dimen2 by-\pagetotal\ifdim\dimen0>\dimen2\newpage\fi\endgroup}
\newcommand{\needspace}[1]{\par\begingroup\skip0=\lastskip\vskip-\skip0\dimen0=#1\relax\advance\dimen0 by\skip0\vskip\dimen0\penalty-100\vskip-\dimen0\vskip\skip0\penalty9999\endgroup}
\usepackage[hidelinks,unicode]{hyperref}
\usepackage[most]{tcolorbox}
\definecolor{ContestBlue}{HTML}{17365D}
\hypersetup{pdftitle={第041套},pdfauthor={竞赛物理教研}}
\setstretch{1.24}
\setlength{\parindent}{2em}
\setlength{\parskip}{0.35em}
\setlength{\emergencystretch}{2em}
\setlist{itemsep=0.35em,topsep=0.4em}
\captionsetup{font=small,labelfont=bf,justification=centering}
\setcounter{secnumdepth}{0}
\pagestyle{fancy}
\fancyhf{}
\fancyhead[L]{\small 第041套}
\fancyfoot[C]{\thepage}
\renewcommand{\headrulewidth}{0.4pt}
\ctexset{
  section={format=\Large\bfseries\color{ContestBlue},beforeskip=1.2em,afterskip=0.7em},
  subsection={format=\large\bfseries\color{ContestBlue},beforeskip=1.0em,afterskip=0.5em},
  subsubsection={format=\normalsize\bfseries,beforeskip=0.8em,afterskip=0.35em}
}
\makeatletter
\newcommand{\tightlist}{\setlength{\itemsep}{0pt}\setlength{\parskip}{0pt}}
\newcommand{\maxwidth}{\ifdim\Gin@nat@width>\linewidth\linewidth\else\Gin@nat@width\fi}
\newcommand{\maxheight}{\ifdim\Gin@nat@height>0.72\textheight 0.72\textheight\else\Gin@nat@height\fi}
\newcommand{\pandocbounded}[1]{\begingroup\setkeys{Gin}{width=\maxwidth,height=\maxheight,keepaspectratio}#1\endgroup}
\makeatother
\begin{document}
\begin{titlepage}
  \centering
  \vspace*{58mm}
  {\Huge\bfseries\color{ContestBlue} 第041套\par}
  \vfill
  \vspace*{28mm}
\end{titlepage}
\section{第 1 题}\label{ux7b2c-1-ux9898}

一根长度为 \(l\)、质量为 \(m\) 的均匀木条，水平放在质量可忽略的木块 A、B
上。A 固定，开始时位于木条左端；B 位于木条右端左侧 \(x_{10}\)
处，\(0<x_{10}<l/2\)。木条与两个木块的动摩擦因数均为 \(\mu_0\)，B
与地面无摩擦。木条和 B 最初静止，给 B 施加水平向左的力 \(F(t)\)，使 B
和木条均作匀加速运动，直至 B 与 A 相遇。求二者的加速度和 \(F(t)\)。

原卷把木条位移 \(x_2(t)\) 写作``向右位移''；以向右为正时，实际运动的
\(x_2\) 为负。下面另设向左位移
\(u=-x_2\)，避免把坐标符号误当成运动方向。

\begin{figure}[H]
\centering
\includegraphics[width=0.9\linewidth,height=\textheight,keepaspectratio,alt={题面模型}]{../figures/question-01.pdf}
\caption{题面模型}
\end{figure}

\newpage

\section{第 2 题}\label{ux7b2c-2-ux9898}

两条相互垂直的光滑细管分别沿水平、竖直方向放置，两质量相等的带电小球
A、B 只能在各自管内运动，电量分别为 \(+q,-q\)。开始 A 静止于距交点 O 为
\(d\) 处，B 位于 O，速度竖直向下，大小
\(v_0=\sqrt{kq^2/(3md)}\)。忽略重力，求两球间距离的变化范围，以及 A
第一次返回初始位置所需时间。

\begin{figure}[H]
\centering
\includegraphics[width=0.9\linewidth,height=\textheight,keepaspectratio,alt={题面模型}]{../figures/question-02.pdf}
\caption{题面模型}
\end{figure}

\newpage

\section{第 3 题}\label{ux7b2c-3-ux9898}

一辆平衡车沿倾角 \(0.8^\circ\) 的斜坡下行。车轮可视为质量
\(2.68\,\mathrm{kg}\)、半径 \(0.174\,\mathrm m\)
的均匀圆柱，其余部分与载荷的总质量为 \(5.12\times10\,\mathrm{kg}\)，取
\(g=9.8\,\mathrm{m/s^2}\)。保持车体姿态、车轮纯滚动。求：(1)
车轮绕自身中轴的转动惯量；(2) 匀速下行时电机的制动力矩；(3) 速度
\(1.18\times10\,\mathrm{m/s}\)、能量回收效率 \(20\%\) 时的发电功率；(4)
关闭电机后车的加速度。

\begin{figure}[H]
\centering
\includegraphics[width=0.9\linewidth,height=\textheight,keepaspectratio,alt={题面模型}]{../figures/question-03.pdf}
\caption{题面模型}
\end{figure}

\newpage

\section{第 4 题}\label{ux7b2c-4-ux9898}

地球半径 \(r=6.37\times10^6\,\mathrm m\)，地面重力加速度
\(g=9.8\,\mathrm{m/s^2}\)。一根均匀柔链的两端由径向支柱固定，端点距地心
\(R=1.27\times10^7\,\mathrm m\)，链的最低点距地心
\(r_1=8.91\times10^6\,\mathrm m\)。端点张力
\(F_1=1.84\times10^7\,\mathrm N\)，最低点张力
\(F_2=8.83\times10^6\,\mathrm N\)。求链在端点与支柱的夹角和链的线密度。

\begin{figure}[H]
\centering
\includegraphics[width=0.9\linewidth,height=\textheight,keepaspectratio,alt={题面模型}]{../figures/question-04.pdf}
\caption{题面模型}
\end{figure}

\newpage

\section{第 5 题}\label{ux7b2c-5-ux9898}

激光干涉仪的两臂长均为 \(L=4\,\mathrm{km}\)，激光波长
\(\lambda=10.6\,\mu\mathrm m\)，每臂往返 \(N=200\)
次。若相位变化可分辨到 \(10^{-8}\)，估计可测相对臂长变化 \(\sigma_0\)
的数量级。月球质量 \(7.35\times10^{22}\,\mathrm{kg}\)、地月距离
\(3.84\times10^8\,\mathrm m\)，以 \(\sigma_M\sim GM/(dc^2)\)
估计月球影响，求对应相位变化的数量级。

设引力波能流密度
\(j=A\sigma^2G^\alpha c^\beta f^\gamma\)，用量纲法确定指数。某事件中
\(29\)、\(36\) 个太阳质量的黑洞并合为 \(62\) 个太阳质量，距离
\(1.3\times10^9\) 光年，释放持续约 \(0.2\,\mathrm s\)，频率约
\(100\,\mathrm{Hz}\)。取
\(A=1\)、\(M_\odot=1.99\times10^{30}\,\mathrm{kg}\)、\(G=6.67\times10^{-11}\)（SI）、\(c=3.0\times10^8\,\mathrm{m/s}\)，各向同性平均估计振幅数量级。原题的波长、多次往返和平均辐射模型按所给条件使用。

\begin{figure}[H]
\centering
\includegraphics[width=0.9\linewidth,height=\textheight,keepaspectratio,alt={题面模型}]{../figures/question-05.pdf}
\caption{题面模型}
\end{figure}

\newpage

\section{第 6 题}\label{ux7b2c-6-ux9898}

长度为 \(l\) 的细狭缝平行于双棱镜的棱，狭缝至棱镜距离为
\(h\)，棱镜至屏距离为 \(L\)。棱镜折射率为 \(n\)，两半的小折射角均为
\(\alpha\)。单色光波长为
\(\lambda\)，狭缝各点非相干发光。求屏上干涉条纹间距和干涉区内的亮条纹数；再把狭缝在自身平面内转过小角
\(\beta\)，求干涉条纹调制度 \(M\)
及其第一次降为零的转角。采用近轴、小棱镜角近似。

\begin{figure}[H]
\centering
\includegraphics[width=0.9\linewidth,height=\textheight,keepaspectratio,alt={题面模型}]{../figures/question-06.pdf}
\caption{题面模型}
\end{figure}

\newpage

\section{第 7 题}\label{ux7b2c-7-ux9898}

第一部分：考虑 \(dy/dx=y^{3/4}\)，\(y\ge0\)。(1)
求含一个待定常数的通解并确定定义域；(2)
给出不能由该正支通解表示的无待定常数特解；(3)
画通解和特解曲线，讨论零点初值的非唯一性及原题关于混沌的说法。

\begin{enumerate}
\def\labelenumi{(\arabic{enumi})}
\setcounter{enumi}{1}
\tightlist
\item
  由质量均为 \(m\)、电荷分别为 \(+q,-q\) 的两小球和长度 \(l\)
  的轻刚杆组成电偶极子，在垂直平面的均匀磁场 \(B\)
  中运动。初始质心静止，杆沿 \(x\) 轴，正电荷在正 \(x\) 侧，初始角速度为
  \(\omega_0>0\)。忽略其他力。建立运动方程，画转角和角动量的相图，确定杆作振动与连续转动的分界，并在振动情形求质心沿初始杆方向的最大位移。
\end{enumerate}

\begin{figure}[H]
\centering
\includegraphics[width=0.9\linewidth,height=\textheight,keepaspectratio,alt={题面模型}]{../figures/question-07.pdf}
\caption{题面模型}
\end{figure}

\newpage

\section{第 8 题}\label{ux7b2c-8-ux9898}

单个质量为 \(m\) 的粒子在势能 \(V(x)=kx^2/2\) 中处于温度 \(T\)
的经典热平衡，总能量为 \(E=mv^2/2+V(x)\)，相空间概率密度正比于
\(e^{-E/(k_BT)}\)。求位置概率密度 \(f(x)\) 和
\(\langle x^2\rangle\)。两个互相独立的同类粒子，求
\(\langle(x_1-x_2)^2\rangle\)。

再令两粒子之间增加势能 \(k(x_1-x_2)^2/2\)，求粒子 1
的边缘位置密度与相对坐标的均方值。可使用高斯积分。

\begin{figure}[H]
\centering
\includegraphics[width=0.9\linewidth,height=\textheight,keepaspectratio,alt={题面模型}]{../figures/question-08.pdf}
\caption{题面模型}
\end{figure}
\end{document}
